\uuid{h8cy}
\titre{Identification des coefficients d'une série entière}
\niveau{L2}
\module{Analyse}
\chapitre{Série entière}
\sousChapitre{Développement en série entière}
\theme{Dérivation terme à terme, série géométrique, parité des coefficients}
\auteur{Maxime Nguyen}
\datecreate{2026-05-19}
\organisation{AMSCC}
\difficulte{2}

\contenu{
	
	\texte{
		On rappelle que, pour tout réel $x$ tel que $|x|<1$,
		\[
		\frac{1}{1-x} = \sum_{n=0}^{+\infty} x^n.
		\]
	}
	
	\begin{enumerate}
		
		\item
		\question{En dérivant cette égalité, montrer que
			\[
			\frac{1}{(1-x)^2} = \sum_{n=0}^{+\infty} (n+1)x^n.
			\]
		}
		\indication{La dérivation terme à terme est licite sur le disque ouvert de convergence.}
		\reponse{
			La série $\sum_{n=0}^{+\infty} x^n$ converge normalement sur tout compact de $(-1,1)$,
			donc on peut dériver terme à terme sur $(-1,1)$.
			On obtient
			\[
			\frac{1}{(1-x)^2} = \left(\frac{1}{1-x}\right)' = \sum_{n=1}^{+\infty} n x^{n-1}
			= \sum_{n=0}^{+\infty} (n+1) x^n.
			\]
		}
		
		\item
		\question{En déduire le développement en série entière de $\dfrac{x}{(1-x)^2}$.
			On donnera explicitement le coefficient de $x^n$ pour tout $n\in\mathbb{N}$.
		}
		\reponse{
			Pour $|x|<1$, en multipliant l'égalité précédente par $x$ :
			\[
			\frac{x}{(1-x)^2} = \sum_{n=0}^{+\infty} (n+1)x^{n+1} = \sum_{n=1}^{+\infty} n\, x^n.
			\]
			Le coefficient de $x^n$ est donc $n$ pour tout $n\in\mathbb{N}$
			(avec la convention que le coefficient de $x^0$ est $0$).
		}
		
		\item
		\question{En utilisant l'égalité $\dfrac{1}{1-x^2} = \dfrac{1}{1-(x^2)}$,
			déterminer le développement en série entière de $\dfrac{1}{1-x^2}$.
			On précisera le coefficient de $x^n$ selon la parité de $n$.
		}
		\reponse{
			Pour $|x|<1$, on substitue $x^2$ à $x$ dans la série géométrique :
			\[
			\frac{1}{1-x^2} = \sum_{n=0}^{+\infty} x^{2n}.
			\]
			Le coefficient de $x^n$ est donc $1$ si $n$ est pair, et $0$ si $n$ est impair.
			Le rayon de convergence est $R=1$.
		}
		
		\item
		\question{Déterminer de même le développement en série entière de $\dfrac{x}{1-x^2}$.
			On précisera le coefficient de $x^n$ selon la parité de $n$.
		}
		\reponse{
			Pour $|x|<1$, en multipliant par $x$ :
			\[
			\frac{x}{1-x^2} = \sum_{n=0}^{+\infty} x^{2n+1}.
			\]
			Le coefficient de $x^n$ est donc $1$ si $n$ est impair, et $0$ si $n$ est pair.
			Le rayon de convergence est $R=1$.
		}
		
		\item
		\question{En déduire le développement en série entière de $\dfrac{1+x}{1-x^2}$.
			Que remarque-t-on ?
		}
		\reponse{
			Pour $|x|<1$, en additionnant les deux développements précédents :
			\[
			\frac{1+x}{1-x^2} = \frac{1}{1-x^2} + \frac{x}{1-x^2}
			= \sum_{n=0}^{+\infty} x^{2n} + \sum_{n=0}^{+\infty} x^{2n+1}
			= \sum_{n=0}^{+\infty} x^n.
			\]
			On remarque que $\dfrac{1+x}{1-x^2} = \dfrac{1+x}{(1-x)(1+x)} = \dfrac{1}{1-x}$
			pour $x \neq -1$, ce qui est cohérent avec la série géométrique de départ.
		}
		
		\item
		\question{Pour chacune des fonctions précédentes, préciser le rayon de convergence
			du développement obtenu.
		}
		\reponse{
			Toutes les séries obtenues ont un rayon de convergence $R = 1$ :
			\begin{itemize}
				\item $\dfrac{1}{(1-x)^2} = \sum_{n=0}^{+\infty}(n+1)x^n$, $R=1$ (la série géométrique dérivée terme à terme conserve le rayon de convergence) ;
				\item $\dfrac{x}{(1-x)^2} = \sum_{n=1}^{+\infty} n\, x^n$, $R=1$ ;
				\item $\dfrac{1}{1-x^2} = \sum_{n=0}^{+\infty} x^{2n}$, $R=1$ ;
				\item $\dfrac{x}{1-x^2} = \sum_{n=0}^{+\infty} x^{2n+1}$, $R=1$ ;
				\item $\dfrac{1+x}{1-x^2} = \sum_{n=0}^{+\infty} x^n$, $R=1$.
			\end{itemize}
		}
		
	\end{enumerate}
}